%% ============================================================================
%%  chapters/chapter1.tex
%%  فصل ۱: کلیات پژوهش
%%  قالب رسمی پایان‌نامه و رساله دانشگاه خوارزمی
%% ============================================================================

\chapter{کلیات پژوهش}
\label{chap:intro}

\section{مقدمه}
\label{sec:ch1_intro}

توسعه روزافزون دانش و فناوری در جهان امروز، ساماندهی، مستندسازی و نشر دقیق یافته‌های علمی را به یکی از ارکان بنیادین پیشرفت جوامع علمی تبدیل کرده است. در این میان، پایان‌نامه‌ها و رساله‌های دانشگاهی، عصاره فعالیت‌های پژوهشی دانشجویان تحصیلات تکمیلی و بازتاب‌دهنده توان علمی دانشگاه‌ها و مراکز تحقیقاتی به شمار می‌روند. ارائه یک اثر علمی مستلزم رعایت دو مؤلفه اساسی است: محتوای غنی علمی و قالب استاندارد نگارشی. هرچند ارزش ذاتی یک رساله در گرو نوآوری‌های نظری، دقت در روش‌شناسی و صحت نتایج تجربی آن است، اما شیوه نمایش، صفحه‌آرایی منظم، یکدستی علائم، و رعایت دقیق دستورالعمل‌های نگارشی دانشگاه نقش تعیین‌کننده‌ای در انتقال صحیح پیام علمی و خوانایی اثر ایفا می‌کند.

در دهه‌های اخیر، سامانه‌های حروف‌چینی دیجیتال تحولی ژرف در تولید مستندات علمی پدید آورده‌اند. سامانه حروف‌چینی تک\index{تک@\lr{TeX}}\LTRfootnote{\lr{TeX}} که توسط دانلد کنوت ابداع گردید و به دنبال آن بسته توسعه‌یافته لاتک\index{لاتک@\lr{LaTeX}}\LTRfootnote{\lr{LaTeX}}، به عنوان ابزارهای معیار در جامعه علمی شناخته شده‌اند. قابلیت‌های منحصر‌به‌فرد این سامانه‌ها در حروف‌چینی دیجیتال\index{حروف‌چینی دیجیتال} فرمول‌های پیچیده ریاضی، ارجاع‌دهی خودکار، مدیریت یکپارچه فهرست‌ها و حفظ ثبات ساختاری متن، موجب تمایز آشکار آن‌ها از واژه‌پردازهای سنتی گردیده است. با این حال، استفاده از این سامانه‌ها در زبان‌هایی با جهت نگارش راست‌به‌چپ\LTRfootnote{\lr{Right-to-Left (RTL)}} نظیر زبان فارسی، همواره با چالش‌های فنی و ساختاری فراوانی روبه‌رو بوده است. ظهور موتور زی‌تک\LTRfootnote{\lr{XeTeX}} و توسعه بسته بومی زی‌پرشین\index{زی‌پرشین@\lr{XePersian}}\LTRfootnote{\lr{XePersian}} زمینه مناسبی را برای حل ریشه‌ای این مسائل و انطباق کامل با ظرافت‌های تایپوگرافی خط فارسی فراهم ساخت.

دانشگاه خوارزمی به عنوان نخستین مؤسسه آموزش عالی مدرن کشور، با تدوین شیوه‌نامه جامع نگارش پایان‌نامه‌ها و رساله‌ها\index{شیوه‌نامه دانشگاه خوارزمی}، استانداردهای دقیقی را برای ساختار، اندازه‌ها، قلم‌ها، حاشیه‌ها و صفحه‌آرایی آثار علمی تعیین کرده است. هدف اصلی این پژوهش، طراحی و پیاده‌سازی یک قالب پایدار، جامع و کاملاً منطبق بر شیوه‌نامه مصوب دانشگاه خوارزمی است تا پژوهشگران بتوانند بدون دغدغه‌های مربوط به تنظیمات ظاهری، بر کیفیت محتوای علمی خود تمرکز نمایند.

\section{بیان مسئله}
\label{sec:ch1_problem}

نگارش متون علمی دانشگاهی به زبان فارسی در بستر سامانه‌های دیجیتال همواره با پیچیدگی‌های چندگانه‌ای روبه‌رو بوده است. از یک سو، واژه‌پردازهای تجاری نظیر مایکروسافت ورد\LTRfootnote{\lr{Microsoft Word}} اگرچه دارای رابط کاربری بصری هستند، اما در مدیریت مستندات حجیم، مدیریت خودکار مراجع دوزبانه، شماره‌گذاری پایدار روابط ریاضی و جدول‌ها، و کنترل دقیق فواصل ساختاری دچار ناپایداری‌های مکرر می‌شوند. به‌هم‌ریختگی جایگاه تصویرها، تغییر ناخواسته قلم‌ها، و عدم تطابق شماره‌گذاری بخش‌ها در اسناد بزرگ از معضلات شناخته‌شده این ابزارهاست.

از سوی دیگر، حروف‌چینی متون فارسی در سامانه لاتک بدون وجود یک قالب دانشگاهی بومی‌سازی‌شده، با دشواری‌های فنی متعددی مواجه است:
\begin{enumerate}
  \item \textbf{تداخل جهت متن دوزبانه:} متون علمی مهندسی و علوم پایه همواره آکنده از اصطلاحات فنی لاتین، اختصارات انگلیسی، و روابط ریاضی در دل پاراگراف‌های فارسی هستند. مدیریت صحیح جهت جریان متن\LTRfootnote{\lr{Bidirectional Text Flow}} به‌گونه‌ای که علائم نگارشی نظیر نقطه، خط‌فاصله، پرانتز و پانویس‌ها در جایگاه درست خود قرار گیرند، نیازمند تنظیم دقیق قواعد دوسویه است.
  \item \textbf{شیوه‌نامه دقیق دانشگاه خوارزمی:} شیوه‌نامه مصوب دانشگاه دارای ضوابط بسیار دقیق در خصوص فواصل حاشیه‌ها (بالا ۳/۵، پایین ۲/۵، راست ۳/۰ و چپ ۳/۰ سانتی‌متر)، فاصله سرصفحه تا لبه کاغذ (۲/۲۵ سانتی‌متر)، اندازه قلم متن اصلی (بی زر ۱۳ با فاصله خطوط ۱/۱۵)، اندازه و فاصله سرعنوان‌ها، و سبک سرصفحه‌ها با ساختار «عنوان \lr{|} شماره صفحه» است که پیاده‌سازی دستی آن توسط دانشجویان بسیار وقت‌گیر و همراه با خطاست.
  \item \textbf{یکپارچگی شماره‌گذاری فارسی:} تمامی عناصر سند شامل شماره صفحه‌ها، شماره بخش‌ها، جدول‌ها، شکل‌ها، نمودارها، روابط ریاضی، الگوریتم‌ها و پانویس‌ها باید از ارقام فارسی تبعیت نمایند.
  \item \textbf{مدیریت مراجع دوزبانه:} درج هم‌زمان مراجع فارسی و لاتین بدون تداخل در سبک کتاب‌شناسی، مرتب‌سازی صحیح و حفظ الگوهای ارجاع‌دهی درون‌متنی، نیازمند سبک‌های اختصاصی بیب‌تک\LTRfootnote{\lr{BibTeX}} است.
\end{enumerate}

بنابراین، عدم وجود یک قالب رسمی، ماژولار و آزموده‌شده برای دانشگاه خوارزمی سبب اتلاف زمان قابل‌توجهی از دانشجویان تحصیلات تکمیلی در مراحل پایانی فراغت از تحصیل و بروز ناهمگونی در انتشارات دانشگاهی می‌گردد.

\section{اهداف پژوهش}
\label{sec:ch1_objectives}

هدف بنیادین این پژوهش، طراحی، پیاده‌سازی و اعتبارسنجی یک بسته و کلاس جامع لاتک برای پایان‌نامه‌ها و رساله‌های دانشگاه خوارزمی است. در راستای تحقق این هدف کلان، اهداف فرعی زیر تعریف شده‌اند:
\begin{itemize}
  \item طراحی کلاس اختصاصی \lr{\texttt{kharazmi-thesis.cls}} با بهره‌گیری از کلاس پایه \lr{\texttt{report}} برای اعمال مستقیم هندسه صفحه، فواصل، سرصفحه‌ها و قلم‌ها.
  \item پیاده‌سازی بسته تکمیلی \lr{\texttt{kharazmi.sty}} جهت تعریف متغیرهای فراداده، صفحات پیش از متن (جلد فارسی، تصویب‌نامه، تعهدنامه، تقدیم، سپاسگزاری، چکیده) و صفحات پس از متن (واژه‌نامه‌ها، کارنامک، فهرست مقالات، جلد انگلیسی).
  \item ایجاد سازوکار مدیریت مراجع دوزبانه با انطباق کامل با موتور زی‌پرشین.
  \item ارائه ۵ فصل نمونه واقعی با پوشش کامل فرمول‌های ریاضی، جداول تخصصی، تصویرها، نمودارها، الگوریتم‌ها و قطعه کدها.
  \item تدوین راهنمای جامع کاربری و مستندات گام‌به‌گام نصب و کامپایل برای استفاده دانشجویان.
\end{itemize}

\section{پرسش‌ها و فرضیه‌های پژوهش}
\label{sec:ch1_questions}

در راستای هدایت دقیق فرآیند طراحی و ارزیابی فنی سامانه، پرسش‌های پژوهشی زیر مطرح می‌گردند:
\begin{itemize}
  \item \textbf{پرسش اول ($Q_1$):} چگونه می‌توان هندسه صفحه، سرصفحه‌ها و فواصل ساختاری تعریف‌شده در شیوه‌نامه دانشگاه خوارزمی را با کمترین وابستگی به بسته‌های متداخل در بستر زی‌لاتک پیاده‌سازی نمود؟
  \item \textbf{پرسش دوم ($Q_2$):} مناسب‌ترین سازوکار برای یکپارچه‌سازی مراجع فارسی و انگلیسی بدون ایجاد خطای جهت متن و نویسه‌های گم‌شده چیست؟
  \item \textbf{پرسش سوم ($Q_3$):} طراحی ماژولار ساختار پایان‌نامه تا چه اندازه در کاهش زمان آماده‌سازی و ارتقای تجربه کاربری دانشجویان مؤثر است؟
\end{itemize}

متناظر با پرسش‌های فوق، فرضیه‌های پژوهشی زیر صورت‌بندی شده‌اند:
\begin{itemize}
  \item \textbf{فرضیه اول ($H_1$):} استفاده از ترکیب کنترل‌شده بسته‌های \lr{\texttt{geometry}}، \lr{\texttt{fancyhdr}} و \lr{\texttt{titlesec}} با تنظیمات سفارشی در کلاس پایه، امکان انطباق میلی‌متری با شیوه‌نامه دانشگاه را فراهم می‌سازد.
  \item \textbf{فرضیه دوم ($H_2$):} بهره‌گیری از سبک کتاب‌شناسی \lr{\texttt{plain-fa}} و فیلد زبانی در فایل \lr{\texttt{.bib}}، پایداری کامل ارجاعات دوزبانه را بدون هشدار خطای نویسه تضمین می‌کند.
  \item \textbf{فرضیه سوم ($H_3$):} تفکیک ساختار قالب به پوشه‌های مجزای صفحات آغازین، فصول و بخش‌های پایانی، نرخ خطای کاربران را به حداقل رسانده و یادگیری الگوها را تسهیل می‌نماید.
\end{itemize}

\section{ضرورت و اهمیت پژوهش}
\label{sec:ch1_importance}

اهمیت این پژوهش را می‌توان از دو بعد دانشگاهی و فنی تبیین نمود. از دیدگاه دانشگاهی، یکدستی و انضباط بصری رساله‌ها نشان‌دهنده هویت حرفه‌ای و انضباط سازمانی دانشگاه خوارزمی است. دیجیتالی‌سازی و انتشار الکترونیکی پایان‌نامه‌ها در کتابخانه‌های دیجیتال نیازمند فایل‌های \lr{PDF} استاندارد با برچسب‌های متنی خوانا، فراداده‌های ساخت‌یافته و پیوندهای فرامتنی پایدار است.

از بعد فنی و کاربردی، وجود این قالب به دانشجویان امکان می‌دهد به‌جای صرف صدها ساعت زمان برای برطرف‌کردن خطاهای ظاهری واژه‌پردازها، با آسودگی خاطر بر نوآوری علمی و تحلیل داده‌های پژوهش خود تمرکز نمایند. همچنین این قالب به عنوان یک الگوی بازتولیدپذیر\LTRfootnote{\lr{Reproducible Research}} برای نگارش مقالات علمی، گزارش‌های پژوهشی و کتب دانشگاهی در آینده قابل استفاده خواهد بود.

\section{ساختار پایان‌نامه}
\label{sec:ch1_organization}

ساختار این پایان‌نامه در قالب ۵ فصل و پیوست‌های مربوطه به شرح زیر سازمان‌دهی شده است:
\begin{itemize}
  \item \textbf{فصل اول (کلیات پژوهش):} به معرفی زمینه تحقیق، بیان مسئله، اهداف، پرسش‌ها، فرضیه‌ها، و ضرورت پژوهش اختصاص دارد.
  \item \textbf{فصل دوم (پیشینه پژوهش و مبانی نظری):} مبانی نظری حروف‌چینی دیجیتال، نظریه ارتباطات، مدل‌های یادگیری عمیق، و پیشینه کارهای پیشین داخلی و بین‌المللی را به تفصیل مرور و مقایسه می‌کند.
  \item \textbf{فصل سوم (روش‌شناسی و مدل پیشنهادی):} به تشریح معماری سامانه، فرمول‌بندی روابط ریاضی، ساختار الگوریتم‌ها و نحوه پیاده‌سازی اجزا می‌پردازد.
  \item \textbf{فصل چهارم (یافته‌ها و تحلیل نتایج):} نتایج تجربی، جداول ارزیابی عملکرد، نمودارهای مقایسه‌ای، شکل‌ها، قطعه کدها و تحلیل‌های کمی و کیفی را ارائه می‌دهد.
  \item \textbf{فصل پنجم (نتیجه‌گیری و پیشنهادات پژوهشی):} جمع‌بندی نهایی، دستاوردها، محدودیت‌های کار، و پیشنهادهایی برای پژوهش‌های آینده را تبیین می‌نماید.
  \item \textbf{بخش‌های پایانی:} شامل پیوست‌ها، واژه‌نامه‌های توصیفی و دوزبانه، فهرست مقالات، کارنامک علمی نویسنده، و مشخصات انگلیسی پایان‌نامه است.
\end{itemize}
